\section{Global completion moments and the drift of an unguided code walk}
\label{sec:completion-moments}

Let $p$ be an odd prime, $N=4p$, and let $G$ be a totally isotropic
subspace of $\mathbb F_p^N$ contained in a self-dual code. Write
$\dim G=N/2-h$, where $h\ge2$. The quotient
$V=G^\perp/G$ is a split nondegenerate orthogonal space of dimension
$2h$. A sign class is normalized by its first coordinate being $+1$.
Every sign vector in $G^\perp$ is isotropic modulo $p$ because its
norm is $N$. Let $b$ count sign classes in $G$ itself, and assign
weight $w_\ell$ to each isotropic projective point $\ell$ of $V$:
it is the number of sign classes whose nonzero quotient lies on
$\ell$. Put
\[
 W=\sum_\ell w_\ell,\qquad S_2=\sum_\ell w_\ell^2,\qquad
 T=\sum_{\substack{\ell\ne k\\\ell\perp k}}w_\ell w_k.
\]
The last sum is ordered. These definitions count actual sign classes,
not arbitrary vectors of the finite-field code.

\begin{theorem}[Exact first and second completion moments]
\label{mom:two}
Choose uniformly among all self-dual codes $C$ containing $G$ and
put $Y=|S(C)|-b$. With $D_1=p^{h-1}+1$ and $D_2=p^{h-2}+1$,
\[
 \mathbb E Y=\frac{W}{D_1},\qquad
 \mathbb E Y^2=\frac{S_2}{D_1}+\frac{T}{D_1D_2}.
\]
If $W>0$, at least one such code satisfies
\[
 |S(C)|\ \ge\
 b+\left\lceil\frac{S_2+T/D_2}{W}\right\rceil.
\]
\end{theorem}
\begin{proof}
The number of maximal totally isotropic subspaces of a split
$2h$-space is
\[
 L_h=2\prod_{j=1}^{h-1}(p^j+1),\qquad L_0=1.
\]
For completeness, the hyperbolic equation is $x\cdot y=0$ with
$x,y\in\mathbb F_p^h$. Its nonzero solutions, modulo scalar, number
$(p^h-1)(p^{h-1}+1)/(p-1)$: when $x\ne0$ there are $p^{h-1}$
choices of $y$, and when $x=0$ there are $p^h-1$ nonzero choices.
Each maximal isotropic subspace has $(p^h-1)/(p-1)$ projective points.
Counting incidences at one point, whose quotient has dimension
$2(h-1)$, gives $L_h=(p^{h-1}+1)L_{h-1}$.
There are two isotropic lines when $h=1$, giving the formula.

A specified isotropic point belongs to a uniformly chosen generator
with probability $L_{h-1}/L_h=1/D_1$. Two distinct orthogonal points
span an isotropic plane and co-occur with probability
$L_{h-2}/L_h=1/(D_1D_2)$. Nonorthogonal points never co-occur.
Expanding $Y=\sum_\ell w_\ell\,1_{\{\ell\subset C/G\}}$ proves both
moment formulas. Finally $Y^2\le(\max Y)Y$, so the maximum is at least
$\mathbb E Y^2/\mathbb E Y$, with the stated integer rounding.
\end{proof}

If $G$ contains five orthogonal real sign rows, the parity theorem
of the tenth pass makes the sign classes of each completion code
pairwise orthogonal over the integers. Thus the condition
\[
 b+\left\lceil(S_2+T/D_2)/W\right\rceil\ge N-3
\]
is sufficient for a Hadamard matrix, by the earlier three-row
completion theorem. It is a count certificate, not a uniform
existence theorem: a proof still has to supply a $G$ and sufficiently
large counts. Finite-polar-space incidence is classical; the
calculation is included to expose precisely the needed input.
The code-neighbor graph is a classical hyperbolic dual polar graph;
see the dual-polar section of Brouwer et al.,
\emph{The smallest eigenvalues of Hamming graphs, Johnson graphs and
other distance-regular graphs with classical parameters},
\url{https://doi.org/10.1016/j.jctb.2018.04.005}.
No priority is claimed for these incidence or spectral ingredients.

\subsection{Exact benchmark where a second moment helps}
Use the first eight stored rows of the twelve-row order-twenty
terminal frame in
\path{../round10/code_traps/search.json}. Their residue span has
dimension eight and contains $b=9$ sign classes. Here $h=2$, and
exact enumeration gives
\[
 W=42,\qquad S_2=96,\qquad T=588,\qquad D_1=6,\quad D_2=2.
\]
The mean total sign count is $9+7=16$, below $N-3=17$.
The second-moment bound instead gives
\[
 \max |S(C)|\ge9+\lceil65/7\rceil=19,
\]
and therefore certifies a full completion. This is a known-order
benchmark, not an unknown-order construction.

Retaining only the first five rows gives $h=5$, $b=5$, and
$W=S_2=1092$, $T=370512$. The same bound certifies only nine total
classes, short of seventeen. A large pool alone does not suffice;
its incidence structure and the retained space both matter.
At order twelve the checker exhausts all eight self-dual completions
of a specified rank-four retained space, verifying both moment
identities directly. Their additional sign counts are all eight.

\subsection{An exact mean drift for random neighbors}
The restricted-neighbor identity of the tenth pass implies a useful
diagnostic. Fix the retained space $G$ and choose each next code
uniformly among self-dual neighbors containing $G$. Define
\[
 d_h=\frac{p^h-1}{p-1},\qquad
 \mu=b+\frac{W}{1+p^{h-1}},\qquad
 \rho_h=\frac{d_h-1-p^{h-1}}{d_h}.
\]
\begin{proposition}[Exact expected sign count along the walk]
For every starting code $C_0\supset G$ and every integer $t\ge0$,
\[
 \mathbb E |S(C_t)|=\mu+\rho_h^t(|S(C_0)|-\mu).
\]
For $h=2$, $\rho_h=0$. For $h\ge3$, $0<\rho_h<1/p$.
\end{proposition}
\begin{proof}
The preceding pass proved
\[
 \sum_{C'\sim C,\ C'\supset G}(|S(C')|-|S(C)|)
 =b+W-|S(C)|-(|S(C)|-b)p^{h-1}.
\]
There are $d_h$ such neighbors. Dividing and rearranging gives
$\mathbb E[|S(C_{t+1})|\mid C_t]
=\mu+\rho_h(|S(C_t)|-\mu)$. Iterate conditional expectations.
For $h\ge2$,
$\rho_h=(p^{h-1}-p)/(p^h-1)$, which gives the stated bounds.
\end{proof}

When a code has more than the global mean number of sign classes,
an unguided neighbor step lowers that number in expectation.
This statement does not rule out improving neighbors or guided
search; the exact order-twenty census in the other branch finds
such neighbors. The formula explains why unrestricted averaging
does not itself prove an improving move. No assertion of universal
positive drift or of all-order Hadamard existence follows.
